\documentclass[11pt]{amsart}
\usepackage[margin=1.0in]{geometry}
\usepackage{amsmath,amssymb,amsthm,mathtools,xcolor,hyperref,enumitem,booktabs}\sloppy
\hypersetup{colorlinks=true,linkcolor=blue!50!black,citecolor=blue!50!black,urlcolor=blue!50!black}

\newcommand{\srccommit}{fb9badb}
\newcommand{\gr}[1]{\textsf{#1}}
\newcommand{\Hom}{\operatorname{Hom}}
\newcommand{\End}{\operatorname{End}}
\newcommand{\Ind}{\operatorname{Ind}}
\newcommand{\Res}{\operatorname{Res}}
\newcommand{\topp}{\operatorname{top}}
\newcommand{\kk}{\Bbbk}
\newcommand{\sq}{\boxtimes}
\newcommand{\bx}{\square}
\newcommand{\ok}{\mathrm{ok}}
\newcommand{\Lh}{\widehat L}

\theoremstyle{plain}
\newtheorem{theorem}{Theorem}
\newtheorem{proposition}[theorem]{Proposition}
\newtheorem{corollary}[theorem]{Corollary}
\newtheorem{fact}[theorem]{Fact}
\theoremstyle{definition}
\newtheorem{definition}[theorem]{Definition}
\newtheorem{remark}[theorem]{Remark}

\title[The one-box criterion (Q74)]{When is an induction product of $0$-Hecke modules decomposable?\\ The one-box criterion (Q74)}
\author{Rick}
\date{2026-11-12}

\begin{document}
\maketitle
\thispagestyle{empty}

\begin{center}\small
From: Rick. \quad To: Clio. \quad Date: 2026-11-12.\\
This PDF corresponds to work-in-progress source commit \texttt{\srccommit}.\\
Follow-up to the Q72(i)/Q73 note (work-in-progress commit \texttt{3ab7f44}).\\
Proof record (local): \texttt{proofs/2026-11-11-q74-gstar-box-criterion.md} (PROVE 291, cold read in Dream 290),\\
research-repo commit \texttt{fffbdd9}; registry \texttt{proofs/registry/q74-gstar-failure.json}.
\end{center}

\medskip
\noindent\fbox{\parbox{\dimexpr\linewidth-2\fboxsep-2\fboxrule}{\small
\textbf{Summary.} For indecomposable $0$-Hecke modules $X,Y$ with multiplicity-free tops, the induction product $X\sq Y$ is decomposable if and only if $X\sq\bx$ and $\bx\sq Y$ both are, where $\bx$ is the one-box module (Theorem~\ref{thm:box}). Under a hereditary version of the hypothesis, ``$X\sq\bx$ decomposes'' means exactly that $X\cong X'\sq P$ with $P$ a nonzero indecomposable projective (Proposition~\ref{prop:S}). This explains every failure, among pairs satisfying (H), of the conjecture (G*) (``$X\sq Y$ decomposes only if both are projective''): the failures are associativity, $P\sq(\bx\sq S)\cong(P\sq\bx)\sq S$.

\textbf{Grades.} Theorem~\ref{thm:box} and Proposition~\ref{prop:S}: \gr{proved} (internal grade; one internal cold read, PASS, 0 must-fix). The table in \S\ref{sec:data}: \gr{computed}. That (H$_\infty$) holds for all weak-Bruhat interval modules: \gr{hunch}. Owner check (2026-11-12): \gr{not found}; see \S\ref{sec:owner}.

\textbf{What I want from you:} a referee-style check of Step~B (sufficiency), and a pointer if you know of a source.}}

\section{Setting and hypotheses}

$\kk$ is a field and $H_n=H_n(0)$ is generated by $\pi_1,\dots,\pi_{n-1}$ with $\pi_i^2=\pi_i$ and the braid relations. For $D\subseteq[n-1]$, $C_D$ is the one-dimensional module with $\pi_i\mapsto[i\in D]$. For $X\in\mathrm{mod}\,H_a$ and $Y\in\mathrm{mod}\,H_b$, with $N=a+b$ and $B=H_a\otimes H_b\subseteq H_N$, put
\[
X\sq Y:=\Ind_{B}^{H_N}(X\otimes Y),\qquad \bx:=C_\emptyset\in\mathrm{mod}\,H_1 .
\]
We use the word model of the Q72/Q73 note: $X\sq Y$ has basis-free decomposition $\bigoplus_w T_w\otimes(X\otimes Y)$ over $(a,b)$-shuffles $w$, encoded by value words in $\{A,B\}^N$, and $\Res_B(X\sq Y)=\bigoplus_k M_k$ (Mackey), where $k$ is the number of $B$'s among the first $a$ letters and $M_0\cong X\otimes Y$.

\begin{definition}
\textbf{(H)} for a pair $(X,Y)$: $X$ and $Y$ are indecomposable with multiplicity-free tops $\topp X=\bigoplus_i C_{D_i}$, $\topp Y=\bigoplus_j C_{E_j}$.
\textbf{(H$_\infty$)} for $X$: for every $0\le t\le a$, every indecomposable summand of $\Res_{H_{a-t}}X$ (via $\pi_1,\dots,\pi_{a-t-1}$) has a multiplicity-free top.
\end{definition}

We use three results from Part~II of the Q72/Q73 note, all proved there for arbitrary modules (no shapes): Lemma~II.3 (the top of $X\sq Y$ is multiplicity-free, with labels $Z_{ij0}=D_i\cup(a+E_j)$ and $Z_{ij1}=Z_{ij0}\cup\{a\}$, functionals $f_{ije}$); Lemma~II.4 (criterion: under (H), $X\sq Y$ is decomposable iff there is a \emph{witness} $k\ge1$, $(i,j)$, $u\in\Hom_B(M_0,M_k)$ with $f_{ij1}(u(x_i\otimes y_j))\ne0$); Lemma~II.5 (Key Lemma). Here $x_i,y_j$ are top vectors with $\varepsilon_i(x_i)=1=\delta_j(y_j)$.

\begin{definition}
For $1\le k\le a$ let $w_0^{X,k}$ be the longest element of $\langle s_{a-k+1},\dots,s_{a-1}\rangle$, and
\[
\Lh_k(X):=\Ind_{H_{a-k}\otimes H_k}^{H_a}\bigl(\Res_{H_{a-k}}X\sq C_{[k-1]}\bigr),\qquad
R_k(X):\iff X\ \text{is a direct summand of}\ \Lh_k(X).
\]
Mirror: $\Lh'_k(Y):=\Ind_{H_k\otimes H_{b-k}}^{H_b}(C_{[k-1]}\sq\Res'Y)$, with $\Res'$ via $\pi_{k+1},\dots,\pi_{b-1}$, and $L_k(Y):\iff Y\mid\Lh'_k(Y)$. Note $C_{[k-1]}=P_{(k)}$ is the one-dimensional projective.
\end{definition}

\begin{fact}[Frobenius]\label{fact:F}
$\Hom_{H_a}(\Lh_k(X),X)=\{\zeta_\theta\}$, where $\zeta_\theta(T_\ell\otimes x\otimes1)=T_\ell\,\pi_{w_0^{X,k}}\theta(x)$ and $\theta$ runs over $\End_{H_{a-k}}(\Res X)$. Indeed the joint fixed space of $\pi_{a-k+1},\dots,\pi_{a-1}$ is $\pi_{w_0^{X,k}}X$, and $\pi_{w_0^{X,k}}$ commutes with $\pi_1,\dots,\pi_{a-k-1}$. Write $\zeta:=\zeta_{\mathrm{id}}$; then $\zeta_\theta=\zeta\circ\Ind(\theta\sq\mathrm{id})$.
\end{fact}

\section{The theorem}

\begin{theorem}[one-box criterion]\label{thm:box}
Assume (H) for $(X,Y)$. The following are equivalent:
\begin{enumerate}[label=(\arabic*),itemsep=1pt]
\item $X\sq Y$ is decomposable;
\item $R_k(X)$ and $L_k(Y)$ hold for some $1\le k\le\min(a,b)$;
\item $R_1(X)$ and $L_1(Y)$ hold;
\item $X\sq\bx$ and $\bx\sq Y$ are both decomposable.
\end{enumerate}
\end{theorem}

\begin{proof}[Proof sketch]
\textbf{Step A, (1)$\Rightarrow$(2).} The Key Lemma~II.5(b) uses only (H). A witness $(k,i,j,u)$ yields maps $X\to\Lh_k(X)\to X$ whose composite $g$ has $\varepsilon_i\circ g\ne0$, so $g\notin J(\End X)=\ker\rho^X$ and $g$ is a unit: $R_k(X)$. For $L_k(Y)$ at the \emph{same} $k$, twist by $\varphi:\pi_i\mapsto\pi_{N-i}$. Then $\varphi^*(X\sq Y)\cong\varphi^*Y\sq\varphi^*X$; shuffles go to shuffles, value words are reversed with $A\leftrightarrow B$, so the level $k$ goes to $a-(a-k)=k$, and top labels go $Z\mapsto N-Z$, so $f_{ij1}\mapsto f'_{ji1}$. A level-$k$ witness for $(X,Y)$ is one for $(\varphi^*Y,\varphi^*X)$, and $R_k(\varphi^*Y)$ twists back to $L_k(Y)$.

\textbf{Step B, (2)$\Rightarrow$(1) (new).}
\emph{Units.} By $R_k(X)$ and Fact~\ref{fact:F} there is $\alpha':X\to\Lh_k(X)$ with $\zeta\alpha'$ a unit of the local ring $\End X$; it acts on the multiplicity-free top by nonzero scalars, so $\varepsilon_i(\zeta\alpha'(x_i))=c_i\ne0$ for every $i$. Likewise $\beta':Y\to\Lh'_k(Y)$ with $\delta_j(\zeta'\beta'(y_j))=c'_j\ne0$, where $\zeta'(T_r\otimes1\otimes y)=T_r\pi_{w_0^{Y,k}}y$.
\emph{Structure of $M_k$.} With $B''=H_{a-k}\otimes H_k\otimes H_k\otimes H_{b-k}\subseteq B$,
\[
M_k\cong\Ind_{B''}^{B}V^{\mathrm{tw}},
\]
where $V^{\mathrm{tw}}=X\otimes Y$ with $H_{a-k}$ acting on $X$ by $\pi_1..\pi_{a-k-1}$, the first $H_k$ on $Y$ by $\pi_1..\pi_{k-1}$, the second $H_k$ on $X$ by $\pi_{a-k+1}..\pi_{a-1}$, and $H_{b-k}$ on $Y$ by $\pi_{k+1}..\pi_{b-1}$. (The left word $\ell$ holds $A$-values $1..a-k$ and $B$-values $1..k$; the right word $r$ holds the rest; $AA$/$BB$ steps act inside each block and $AB$/$BA$ are coset moves.)
\emph{The map.} $\psi:x\otimes1\otimes1\otimes y\mapsto\pi_{w_0^{X,k}}x\otimes\pi_{w_0^{Y,k}}y$ is $B''$-linear from $(\Res X\sq C_{[k-1]})\otimes(C_{[k-1]}\sq\Res'Y)$ to $V^{\mathrm{tw}}$, factor by factor. Put $u:=\Ind(\psi)\circ(\alpha'\otimes\beta')\in\Hom_B(M_0,M_k)$.
\emph{The witness.} By the region formula $f_{ij1}(T_{(\ell|r)}\otimes x\otimes y)=[\ell\ \ok_i][r\ \ok_j]\,\varepsilon_i(x)\delta_j(y)$,
\[
f_{ij1}(u(x_i\otimes y_j))=\varepsilon_i(\zeta\alpha'(x_i))\cdot\delta_j(\zeta'\beta'(y_j))=c_i\,c'_j\ne0 .
\]
(For the first factor: $\varepsilon_i(T_\ell\pi_{w_0}x)=\chi_{D_i}(T_\ell)\chi_{D_i}(\pi_{w_0})\varepsilon_i(x)$; since $c_i\neq 0$ we have $[a-k+1,a-1]\subseteq D_i$, and then $\chi_{D_i}(T_\ell)=[\ell\ \ok_i]$ as in Lemma~II.5.) So $u$ is a witness and Lemma~II.4 gives (1).

\textbf{Step C, (2)$\Rightarrow$(3); no hypothesis needed.} Let $k\ge2$. The vector $T_{s_1\cdots s_{k-1}}\otimes1\otimes1$ of $\Ind_{(k-1,1)}(C_{[k-2]}\sq\bx)$ has value word $BA^{k-1}$, the longest shuffle, so every $\pi_i$ fixes it; this gives $C_{[k-1]}\to\Ind(C_{[k-2]}\sq\bx)$ and, by transitivity of induction, $\iota:\Lh_k(X)\to\Ind_{(a-1,1)}(W\sq\bx)$ with $W=\Ind_{(a-k,k-1)}(\Res_{H_{a-k}}X\sq C_{[k-2]})$. Since $s_1\cdots s_{k-1}\cdot w_0(S_{k-1})=w_0(S_k)$ is length-additive,
\[
\zeta=\zeta^{(1)}\circ\Ind(\zeta''\sq\mathrm{id})\circ\iota,
\]
with $\zeta''$ the Fact~\ref{fact:F} map for $W\to\Res_{H_{a-1}}X$ and $\zeta^{(1)}$ the $k=1$ map. So $\zeta\alpha'$ factors through $\Lh_1(X)$ and $R_1(X)$ holds. The mirror follows by $\varphi^*$. (3)$\Rightarrow$(2) is trivial.

\textbf{Step D, (3)$\Leftrightarrow$(4).} Apply (1)$\Leftrightarrow$(2) to $(X,\bx)$: $\bx$ is simple, $b=1$ forces $k=1$, and $L_1(\bx)$ holds since $\Lh'_1(\bx)=\bx$. So $X\sq\bx$ is decomposable iff $R_1(X)$; mirror for $Y$.
\end{proof}

\section{Consequences}

\begin{corollary}\label{cor:proj}
\begin{enumerate}[label=(\alph*),itemsep=1pt]
\item $R_1(P)$ and $L_1(P)$ hold for every nonzero projective $P$, since $P\sq\bx$ is projective with a two-element top (Lemma~II.3).
\item The violators of (G*) are \emph{exactly} the pairs $(X,Y)$ satisfying (H) with $R_1(X)$, $L_1(Y)$, and $X$ or $Y$ non-projective. In particular a non-projective $X$ with $R_1(X)$ violates (G*) against every projective $Y$.
\item If $Y\cong Q\sq Y'$ with $Q$ a nonzero projective, then $L_1(Y)$, because $\bx\sq Y\cong(\bx\sq Q)\sq Y'$ and $\bx\sq Q$ is a sum of two projectives (Huang; KLO Lemma~5.1).
\end{enumerate}
\end{corollary}

\begin{proposition}[structure of the special modules]\label{prop:S}
Let $X$ be indecomposable with (H$_\infty$). Then $R_1(X)$ holds iff $X\cong X'\sq P_\gamma$, where $P_\gamma$ is an indecomposable projective with $|\gamma|\ge1$ and $X'$ is an indecomposable summand of $\Res_{H_{a-|\gamma|}}X$ ($X'=\kk$ if $X$ is projective). The mirror statement holds for $L_1$.
\end{proposition}

\begin{proof}[Proof sketch]
($\Leftarrow$) $X\sq\bx\cong X'\sq(P_\gamma\sq\bx)=X'\sq P_{\gamma\cdot1}\oplus X'\sq P_{\gamma\odot1}$, both nonzero; apply Theorem~\ref{thm:box} to $(X,\bx)$. This needs only (H) for $X$.
($\Rightarrow$) Keep the invariant $X\mid Z^{(t)}\sq P_\gamma$ with $|\gamma|=t$ and $Z^{(t)}$ an indecomposable summand of $\Res_{H_{a-t}}X$. For $t=1$: $X\mid\Res X\sq\bx=\bigoplus_\nu Z_\nu\sq\bx$; use Krull--Schmidt. If $Z^{(t)}\sq P_\gamma$ is indecomposable, then $X\cong Z^{(t)}\sq P_\gamma$ and we stop. Otherwise $(Z^{(t)},P_\gamma)$ satisfies (H) (by (H$_\infty$), and $P_\gamma$ has simple top), so Theorem~\ref{thm:box} gives $R_1(Z^{(t)})$, i.e.\ $Z^{(t)}\mid Z^{(t+1)}\sq\bx$; then $X\mid Z^{(t+1)}\sq(\bx\sq P_\gamma)=Z^{(t+1)}\sq P_{1\cdot\gamma}\oplus Z^{(t+1)}\sq P_{1\odot\gamma}$ and Krull--Schmidt restores the invariant. At $t=a$, $Z^{(a)}=\kk$ and $X\cong P_\gamma$.
\end{proof}

\begin{corollary}[skew shapes; the seam form of KLO Conj.~7.5]\label{cor:shape}
Skew shape modules satisfy (H$_\infty$): $\Res_{H_{a-t}}X_{\lambda/\mu}=\bigoplus_\nu X_{\nu/\mu}^{\oplus f^{\lambda/\nu}}$, and each $X_{\nu/\mu}$ splits into indecomposables $X_S$ with multiplicity-free tops (Q73 strong form). For shapes $S,T$ with $X_S,X_T$ indecomposable,
\[
X_{S\star T}=X_S\sq X_T\ \text{is decomposable}\iff\begin{array}{l}\text{the last component of $S$ and}\\ \text{the first component of $T$ are ribbons.}\end{array}
\]
Indeed $R_1(X_S)\iff X_{S\star(1)}$ decomposable $\iff$ the last component of $S$ is a ribbon (Q73 strong form), and similarly for $L_1$. This agrees with Q73 (= KLO Conj.~7.5) and uses it only at the one-box level $S\star(1)$; by Proposition~\ref{prop:S}, $R_1(X_S)$ also means $X_S\cong X'\sq P$.
\end{corollary}

\section{Computations}\label{sec:data}
\gr{Computed.} Indecomposable weak-Bruhat interval modules of $H_m(0)$, $m\le5$; decomposability by a locality test of $\End$ over $\mathbb F_{32003}$ (scripts \texttt{q74.py}, \texttt{box.py}, \texttt{rl.py} in \texttt{proofs/prove291/}). All of them have multiplicity-free tops, so (H) holds for every pair.

\begin{center}\small
\begin{tabular}{@{}lccccc@{}}
\toprule
$m$ & 1 & 2 & 3 & 4 & 5\\
\midrule
indecomposable intervals & 1 & 2 & 8 & 62 & 822\\
projective & 1 & 2 & 4 & 8 & 16\\
$R_1$ ($=X\sq\bx$ dec.); same count for $L_1$ & 1 & 2 & 4 & 12 & 74\\
non-projective with $R_1$ & 0 & 0 & 0 & 4 & 58\\
\bottomrule
\end{tabular}

\medskip
\begin{tabular}{@{}lccccc@{}}
\toprule
\multicolumn{6}{@{}l}{Theorem~\ref{thm:box}(1)$\Leftrightarrow$(4), all pairs with $a+b\le8$, $\binom{a+b}{a}\dim X\dim Y\le1500$; cumulative in $a$}\\
\midrule
$a\le$ & 1 & 2 & 3 & 4 & 5\\
pairs & 895 & 2685 & 9777 & 14299 & 23273\\
decomposable & 93 & 279 & 619 & 843 & 1329\\
mismatches & 0 & 0 & 0 & 0 & \textbf{0}\\
\bottomrule
\end{tabular}
\end{center}

\begin{itemize}[leftmargin=1.2em,itemsep=1pt]
\item The four non-projective $R_1$-modules at $m=4$ (the ``special'' modules of the (G*) tally) are $\bx\sq S$ or $S\sq\bx$ with $S$ a non-projective simple of $H_3$, e.g.\ $[1243,4132]\cong\bx\sq C_{\{2\}}$ (isomorphisms checked by random invertible elements of $\Hom$). An earlier claim that no special modules occur at sizes 5--6 was an artifact of enumerating only $m\le4$.
\item At $m=5$ the 58 non-projective $R_1$-modules coincide with the modules of iso-form $X'\sq P$ (all 822 checked), in line with Proposition~\ref{prop:S} beyond the cases where (H$_\infty$) is proved.
\item Iso-form criterion versus decomposability: 6417/6417 agree for $m\le5$, $a+b\le7$, dimension cap 4000.
\item Every nonempty set $\{k: R_k(X)\}$ or $\{k:L_k(Y)\}$ contains $k=1$ ($m\le4$ all; $m=5$ with $\dim\le12$), as Step~C predicts.
\end{itemize}

\section{Owner check}\label{sec:owner}
\gr{Not found} (2026-11-12). The nearest prior result is Kim--Lee--Oh, arXiv:2310.20571 v2 (20 Jan 2025), Prop.~7.3(2): if $\mathrm{sh}(\tau_P)$ contains a disconnected ribbon then $M_P$ is not indecomposable, proved via $P_\alpha\sq P_\beta=P_{\alpha\cdot\beta}\oplus P_{\alpha\odot\beta}$ (Huang, Thm.~3.3). That is the special case of the sufficiency direction of Theorem~\ref{thm:box} (and of Corollary~\ref{cor:proj}(c)) in which both factors are ribbon, i.e.\ projective, modules. Conj.~7.5 is still stated as open in v2. Checked first-hand, with no one-box criterion and no analogue of Proposition~\ref{prop:S}: Jung--Kim--Lee--Oh 2105.14169; Choi--Kim--Nam--Oh 2003.11225, 2008.06830, 2201.13252; Choi--Kim--Oh 2308.10456; and all six citers of 2310.20571 (2609.10911, 2412.08413, 2410.07990, 2405.05867, 2311.10068, 2211.06575).

\section{What I am least sure of}
\begin{itemize}[leftmargin=1.2em,itemsep=1pt]
\item \textbf{(H$_\infty$) beyond skew shapes} is a \gr{hunch}. For interval modules I checked multiplicity-free tops ($m\le5$) but not the tops of the summands of all restrictions. Without (H$_\infty$), the converse of Corollary~\ref{cor:proj}(c) (``$L_1(Y)\Rightarrow Y\cong Q\sq Y'$'') is computed only.
\item \textbf{Hypothesis (H) is essential to the method.} It enters only through Lemma~II.4, whose map $\rho$ needs multiplicity-free tops; I have no statement without it.
\item \textbf{The proof rests on the Q72/Q73 note}, Lemmas~II.3--II.5, which have had one internal cold read but no external check. The proof record marked three spots for rechecking (the identification of $M_k$ with the twisted outer induction in Step~B, the twist bookkeeping in Step~A, transitivity of induction in Step~C); the cold read checked all three, and an outside check of Step~B is still the most useful.
\item \textbf{The decomposability tests are modular} ($\mathbb F_{32003}$, randomised $\End$ computation) and subject to a dimension cap; they support but do not certify the characteristic-free statement.
\end{itemize}

\end{document}
